% Standalone section for inclusion in the main paper.
% Requires amsmath, amssymb, amsthm; uses the main paper's theorem/lemma envs.
\section{The reverse consistency implication}
\label{sec:reverse-consistency}

The published PA196 endpoint explicitly supplies
\[
 \mathrm{PA}\vdash \operatorname{Nonhalt}(M_{196})
                   \longrightarrow \operatorname{Con}(\mathrm{PA}).
\]
Its external soundness and nonhalting theorems do not, merely by being
formalized in Lean, assert the opposite implication inside PA.
We give a syntactic interpretation of the command calculus and a conditional
transfer theorem for the physical compilation.  The calculus interpretation
below is an ordinary mathematical proof.  Its application to the exact
arithmetized tape predicate still has the explicit implementation and binding
obligations listed at the end of this section.  No new closed Lean theorem
of the internal biconditional is claimed.

\subsection{An arithmetical interpretation of the exact calculus}

Write $B(u,v)$ for the arithmetical predicate asserting that the $v$th
binary digit of $u$ is one.  Thus the first argument is the container,
as in the convention for $\mathsf M$ in PA196.  For convenience work
first in a fixed definitional extension $\mathrm{PA}^{+}$ for the
primitive recursive functions needed to express binary digits and
exponentiation.  Fix their PA proofs of totality and the usual
primitive recursive elimination of this finite definitional extension
from proofs.  In particular,
\begin{equation}\label{eq:bit-bound}
 \mathrm{PA}^{+}\vdash B(u,v)\longrightarrow v<u,
 \qquad
 \mathrm{PA}^{+}\vdash \neg B(u,u).
\end{equation}
The first implication follows from $2^v\leq u$ and $v<2^v$.

The translation $T$ sends truth to $0=0$ and commutes with implication
and universal quantification.  On atoms it is
\[
 T(\mathsf M(a,b))=B(x_a,x_b),\qquad
 T(\mathsf E(a,b))=(x_a=x_b),\qquad
 T(\mathsf U(a,F))=\forall x_a\,T(F).
\]
The numerical decoder reads
$\mathsf U(a,F)=4p(F,a)+2$ with its transposed arguments.
All its recursion is bounded by the input code.  The one degenerate
implication constructor is $\mathsf I(0,0)=0$; here $T(0)=0=0$ is
provably equivalent to $T(0)\to T(0)$, and a fixed logical proof
bridges the syntactic mismatch.  Other implication constructors commute
literally with $T$.

Every variable label in a formula of code $c$ is less than $c$.
This is the syntactic bound proved in
\texttt{PA196/Logic/Decode.lean}; in particular $x_c$ does not occur
in $T(c)$.  No assertion that $T(c)$ is true is involved.

\subsection{Uniform finite-comprehension proofs}

\begin{lemma}[A proof-producing comprehension schema]
For every arithmetical formula $\psi(z,\vec u)$, with auxiliary variables
chosen fresh, $\mathrm{PA}^{+}$ proves
\begin{equation}\label{eq:finite-comprehension}
 \forall n\,\exists s<2^n\,
 \forall z<n\bigl(B(s,z)\leftrightarrow\psi(z,\vec u)\bigr).
\end{equation}
There is a primitive recursive function of the code of $\psi$ that
produces such a proof, and PA proves the syntactic correctness of
this proof-producing function.
\end{lemma}

\begin{proof}
Use the displayed existential formula as the induction formula in $n$.
At $n=0$ take $s=0$.  At the step, take $s+2^n$ if
$\psi(n,\vec u)$ and take $s$ if $\neg\psi(n,\vec u)$.
The binary-digit update laws prove the required statement at $n+1$;
the old witness satisfies $s<2^n$, so adding $2^n$ introduces exactly
the new digit.  This is ordinary classical reasoning in PA.

The proof consists of a fixed finite logical template, the fixed binary
arithmetic lemmas, and one induction-axiom instance built from $\psi$.
Variable renaming, substitution, formation of that instance, and proof
concatenation are primitive recursive syntax operations.  Verifying
that the constructed instance satisfies the PA induction schema is a
syntactic calculation.  Consequently PA verifies the output as a
proof uniformly in the code of $\psi$; no uniform arithmetical truth
predicate is used.
\end{proof}

Applying the lemma to
$\bigl(B(A,z)\lor z=b\bigr)\land\psi(z,\vec u)$,
with cutoff $A+b+1$, gives
\begin{equation}\label{eq:pointed-bits}
 \exists s\,\forall z\left[
 B(s,z)\leftrightarrow
 \bigl((B(A,z)\lor z=b)\land\psi(z,\vec u)\bigr)\right].
\end{equation}
Equation~\eqref{eq:bit-bound} shows that every eligible $z$ lies below
the cutoff, while a witness less than $2^{A+b+1}$ has no digits beyond
it.  This is the finite pointed-separation schema needed by the gate.

\subsection{Translation of the gate and the logical heads}

The six PA196 heads and the restricted propositional head are valid
first-order or equality schemata after translation.  Finite case
distinctions on coinciding variable labels give uniform proof
constructors.  The domination restriction on the fourth argument of
the propositional head restricts the available instances, but plays no
role in their validity.

For the gate, keep the notation of the PA196 paper:
\[
 \Delta=\mathsf E(a,b)\to\mathsf U(c,\mathsf E(a,b)),
 \quad V=\mathsf M(a,b)\to\mathsf U(c,\mathsf M(a,b)),
\]
\[
 P=c\to\mathsf M(c,b),\quad
 S=\mathsf U\bigl(b,(\mathsf M(a,b)\to P)
                         \land(\mathsf M(c,b)\to c)\bigr),
 \quad Q=\mathsf U(c,S\to\neg P).
\]
The returned formula is $\neg(\Delta\to(V\to Q))$ under
$c\ne a,b$, and truth otherwise.

If the guard fails, the translation is immediately provable.  If it
holds, $x_c$ is absent from the translated equality and membership
atoms, so $T(\Delta)$ and $T(V)$ have elementary vacuity proofs.
Put $\varphi=T(c)$; $x_c$ is also absent from $\varphi$ by the code
bound.  For $a\ne b$, equation~\eqref{eq:pointed-bits}, with
$\psi(z)=\varphi[z/x_b]$, constructs a witness $s$ satisfying
\[
 \forall x_b\,
 \bigl[(B(x_a,x_b)\to(\varphi\to B(s,x_b)))
       \land(B(s,x_b)\to\varphi)\bigr]
\]
and also $\varphi\to B(s,x_b)$ at the distinguished original value
of $x_b$.  Thus $\exists x_c(T(S)\land T(P))$ has a PA proof.

The collision $a=b$ deserves a separate check: the first membership
atom inside $S$ then becomes $B(x_b,x_b)$, which PA refutes.  Use
pointed comprehension on the singleton distinguished point; it proves
the second conjunct and the point clause, while the first conjunct
is vacuous.  This handles the collision without incorrectly treating
$x_a$ as an unchanged parameter under the $x_b$ binder.

Classical first-order reasoning now gives $\neg T(Q)$ and hence
$T(\neg(\Delta\to(V\to Q)))$.  Every step is part of a uniform
finite proof template.  This proves the following syntactic statement.

\begin{lemma}[Refutation translation]\label{lem:reverse-proof}
Let $\operatorname{Prf}_{C}(p,f)$ be a primitive recursive proof
predicate for exactly the PA196 numeric calculus.  There is a
primitive recursive function $F$ for which PA proves
\[
 \operatorname{Prf}_{C}(p,1)
       \longrightarrow
 \operatorname{Prf}_{\mathrm{PA}}(F(p),\ulcorner\bot\urcorner).
\]
\end{lemma}

\begin{proof}
Primitive recursion over the finite derivation translates each axiom
using the preceding constructors, each modus ponens by proof
concatenation, and each generalization by the corresponding PA rule.
Use universal closures at intermediate open conclusions.  The
code-zero constructor collision is bridged by the fixed equivalence
already described.  The result is a proof in $\mathrm{PA}^{+}$ of
the translated conclusion.  At code $1$ this is $B(x_0,x_0)$, whose
negation has the fixed proof in~\eqref{eq:bit-bound}.  Append that proof
and eliminate the definitional extension.

Correctness is proved by induction on derivation syntax, using only
syntactic checking of the produced proofs.  This is neither a
reflection principle for PA nor an induction on the truth of
arbitrary PA formulas.
\end{proof}

\subsection{From a finite halting trace to a refutation}

The command invariant in
\texttt{PA196/Logic/CommandSoundness.lean} can be made explicitly
proof-producing.  Retain a derivation code for each live formula $t,h$
and a sequence of derivation codes parallel to the numeric stack.
Initially all formulas are truth.  Stack underflow returns truth with
its fixed proof.  A successful MP appends an application of modus
ponens; a failed MP supplies truth.  AX supplies its six heads,
restricted propositional head, generalization of the old $h$, saved
old top, and guarded gate or truth.  These are primitive recursive
updates of finite data.

When Shift halts, $t+h=1$.  Since these are natural numbers, one of
$t,h$ is exactly $1$, and its retained code is a derivation of the
contradiction.  Notice that this extraction requires no assertion
that a derivation of $1$ cannot exist.

For the tape machine, the finite startup and total positive-length
instruction contracts permit extraction along a finite trace.
Determinism identifies each completed instruction boundary with the
source execution.  A tape trace of length $N$ has at most $N$
completed source instructions, since every instruction contract has
positive length.  No nonhalting engine macro has an internal halting
exit: its totality and designated exit are part of its contract.
Thus a halting trace reaches the source Shift test, and the preceding
annotated invariant returns a calculus refutation.  Bounded traversal
of the trace and the primitive recursive syntax updates make the
extractor primitive recursive.

This gives the following useful transfer criterion.  The published
development contains substantial physical and internal-compilation
infrastructure for its application.  A modified table must satisfy the same
stated hypotheses; the existence of finite test runs does not supply them.

\begin{theorem}[Internal equivalence from certified compilation]
Suppose a binary table $M$ has a PA-verifiable finite startup and
total positive-length instruction contracts for the PA196 command
search, including its unique halting test.  Suppose also that the
PA196 compiler from a coded PA refutation to a finite halting
execution is retained.  Then
\[
 \mathrm{PA}\vdash
 \operatorname{Nonhalt}(M)\leftrightarrow
 \operatorname{Con}(\mathrm{PA}).
\]
\end{theorem}

\begin{proof}
The forward implication is supplied by the refutation-to-execution
compiler, contraposed.  For the reverse implication, PA verifies the
bounded extraction from a halting trace to a calculus refutation and
then applies Lemma~\ref{lem:reverse-proof}.  Hence PA proves that
halting yields a PA refutation.  Contraposition is the desired reverse
implication.
\end{proof}

The theorem identifies the exact extra obligations behind an internal
biconditional.  External soundness, finite testing of selected runs,
and a theorem of two-sided independence are each insufficient by
themselves.  Conversely, no uniform truth predicate and no PA
reflection assumption are needed once these syntactic compilations
and total physical contracts are supplied.

\paragraph{Remaining artifact-level obligations.}
The following concrete work is not supplied by this section and must not
be inferred from the existence proof of the calculus compiler:
\begin{enumerate}
\item Fix the certificate encodings and connect
$\operatorname{Prf}_{C}$ and $\operatorname{Prf}_{\mathrm{PA}}$ above to the
exact checker and consistency sentence used in the published arithmetic
statement.  Record and verify the encoding conversions.
\item Implement the proof-producing templates and their arithmetized
correctness proofs, including the code-zero implication collision, fresh
variables in comprehension, and the coincident-label gate case.  The
mathematical template construction above does not constitute a compiled
Lean implementation of these transformations.
\item Bind the reverse trace extractor to the literal machine's exact
arithmetic trace predicate.  Verify the startup, every relevant positive
instruction contract, termination of its bounded traversal, and extraction
of the Shift certificate within the chosen arithmetization.
\item For a reordered AX body, verify its physical compilation and update
the axiom-selection ranks in the forward certificate compiler.  For push
order $(h_1,h_2,h_3,h_5,h_6,h_4)$ the ranks from the top are respectively
$(6,5,4,3,2,1)$ in that displayed order; equivalently,
$h_4$ has rank $1$, $h_5$ rank $3$, and $h_6$ rank $2$.
The propositional head, generalization and saved top retain ranks $7,8,9$.
\end{enumerate}
The fixed permutation preserves the exact set of available axiom
instances.  It changes which head a particular command word selects;
the appropriate equivalence is through the adjusted proof compiler,
not equality of the literal command executions.
